\section{Lý do chọn đề tài}

Trong kỷ nguyên của Cách mạng Công nghiệp 4.0, \ac{iot} đang phát triển nhanh và được ứng dụng rộng rãi trong nhiều lĩnh vực như nhà thông minh, y tế, giao thông, công nghiệp thông minh,... Sự bùng nổ về số lượng thiết bị \ac{iot} dẫn đến sự gia tăng đáng kể về lưu lượng và loại hình dữ liệu truyền trong mạng. Tuy nhiên, các thiết bị \ac{iot} thường có khả năng xử lý hạn chế và thiếu các cơ chế bảo mật mạnh, khiến chúng trở thành mục tiêu dễ bị tấn công bởi tin tặc.

Để đối phó với các mối đe dọa an ninh mạng, các hệ thống phát hiện xâm nhập
(\ac{ids}) đang được nghiên cứu và phát triển rộng rãi. Đặc biệt, trong những năm gần đây, việc kết hợp học máy (\ac{ml}) với các nền tảng xử lý dữ liệu lớn giúp hệ thống có khả năng phát hiện các tấn công mới (zero-day) và thích nghi với môi trường mạng thay đổi liên tục.
Tuy nhiên, phần lớn các nghiên cứu hiện nay tập trung vào việc xây dựng mô
hình học máy trên tập dữ liệu tĩnh, dung lượng nhỏ và huấn luyện ngoại tuyến (offline).

Trong khi đó, mạng \ac{iot} hiện đại sinh ra lượng dữ liệu rất lớn, liên tục theo thời gian thực đòi hỏi các hệ thống \ac{ids} có khả năng xử lý dữ liệu lớn (\ac{bd}) hiệu quả, có thể hoạt động trong môi trường phân tán và thời gian thực.

Apache Spark là một nền tảng xử lý dữ liệu lớn, với khả năng xử lý
song song, phân tán và dễ dàng tích hợp với thư viện học máy Spark MLlib là một giải pháp tiềm năng để xây dựng hệ thống \ac{ids} hiện đại đáp ứng yêu cầu về quy mô và thời gian thực trong môi trường \ac{iot}.

\section{Các nghiên cứu liên quan}

Trong những năm gần đây, các hệ thống phát hiện xâm nhập (Intrusion Detection System -- \ac{ids}) dựa trên học máy và học sâu đã thu hút nhiều sự quan tâm của cộng đồng nghiên cứu nhằm nâng cao độ chính xác và khả năng thích ứng trước các hình thức tấn công mạng ngày càng tinh vi. Nhiều công trình tập trung khai thác các kiến trúc học sâu như autoencoder, mạng tích chập (\ac{cnn}) và mạng hồi tiếp (\ac{rnn}, \ac{lstm}) để tự động trích xuất đặc trưng và phát hiện bất thường trong lưu lượng mạng. Tuy nhiên, phần lớn các nghiên cứu hiện nay vẫn được triển khai trong môi trường đơn máy, xử lý dữ liệu tĩnh với quy mô tương đối nhỏ, và chưa tích hợp các nền tảng xử lý dữ liệu lớn như Apache Spark hay Hadoop nhằm đáp ứng yêu cầu về khả năng mở rộng và xử lý thời gian thực trong các hệ thống mạng hiện đại.

Một số nghiên cứu đã đề xuất các mô hình \ac{ids} dựa trên autoencoder kết hợp với \ac{cnn} và \ac{lstm} nhằm nâng cao khả năng học đặc trưng và phát hiện xâm nhập mạng. Các mô hình này được đánh giá trên các tập dữ liệu chuẩn như NSL-KDD, CIC-DDoS2019 và CICIDS2017, cho thấy hiệu quả cao trong việc phát hiện bất thường. Tuy nhiên, quá trình huấn luyện và đánh giá vẫn chủ yếu được thực hiện trong môi trường cục bộ và theo cơ chế xử lý batch \cite{singh2022,sun2020}. Tương tự, các mô hình Autoencoder-\ac{lstm} được tối ưu hóa cho bài toán phát hiện bất thường trong mạng cũng đạt được kết quả khả quan, song vẫn triển khai trên nền tảng đơn nút (single-node) và chưa khai thác được khả năng mở rộng của các hệ thống Big Data \cite{dash2025,bamber2024}. Gần đây, một số nghiên cứu đã triển khai \ac{ids} học máy trực tiếp trên phần cứng biên NVIDIA Jetson \cite{hasan2025iiot,baalaji2025ensemble,asghar2026hedid}, nhưng vẫn ở quy mô đơn nút và chưa tích hợp cơ chế xử lý luồng phân tán.

Bên cạnh các phương pháp học sâu, các kỹ thuật học máy truyền thống như Support Vector Machine (\ac{svm}), k-Nearest Neighbors (\ac{knn}) và Decision Tree vẫn được sử dụng rộng rãi trong các hệ thống \ac{ids}. Một số nghiên cứu đã tiến hành so sánh hiệu năng của các mô hình này trên các bộ dữ liệu phổ biến như NSL-KDD và UNSW-NB15, qua đó chỉ ra rằng các phương pháp học máy truyền thống có tiềm năng ứng dụng thực tiễn cao. Tuy nhiên, các mô hình này chủ yếu hoạt động ở chế độ offline, xử lý dữ liệu theo từng đợt và chưa đáp ứng được yêu cầu về xử lý dữ liệu liên tục trong môi trường mạng động \cite{wijden2021}. Ngoài ra, một số công trình đề xuất các mô hình \ac{ids} hai giai đoạn kết hợp giữa học sâu và học máy, chẳng hạn như kết hợp \ac{lstm} và Random Forest, cho kết quả tốt trên các bộ dữ liệu lớn, nhưng toàn bộ thực nghiệm vẫn được thực hiện trên một máy đơn \cite{allail2023}.

Ngoài các kiến trúc nêu trên, nhiều mô hình học sâu khác như \ac{cnn}, \ac{rnn} và các mô hình lai \ac{cnn}-\ac{lstm} cũng đã được nghiên cứu và đánh giá cho bài toán phát hiện xâm nhập mạng. Các nghiên cứu này cho thấy khả năng cải thiện đáng kể độ chính xác so với các phương pháp truyền thống khi thử nghiệm trên các bộ dữ liệu như CICIDS2018. Tuy nhiên, phần lớn các mô hình vẫn thiếu khả năng triển khai trong môi trường thời gian thực và chưa tận dụng được sức mạnh của các nền tảng xử lý dữ liệu lớn \cite{sensors2023}.

Trong bối cảnh Internet of Vehicles (IoV), một số nghiên cứu đã đề xuất các hệ thống \ac{ids} dựa trên học máy và học sâu nhằm bảo vệ các mạng phương tiện thông minh. Các phương pháp này bao gồm việc sử dụng các mô hình dựa trên cây (tree-based), các hệ thống phát hiện xâm nhập đa tầng, cũng như các mô hình ensemble để cải thiện độ chính xác phát hiện tấn công. Mặc dù đạt được kết quả khả quan trên các bộ dữ liệu thử nghiệm, các hệ thống này vẫn chủ yếu được huấn luyện và đánh giá trong môi trường ngoại tuyến, xử lý dữ liệu theo batch và chưa tích hợp các nền tảng xử lý dữ liệu lớn hoặc cơ chế xử lý dữ liệu theo luồng \cite{yang2019,yang2022a,yang2022b}.

Bên cạnh đó, một số hướng tiếp cận gần đây đã nghiên cứu việc áp dụng học bán giám sát và học liên kết (federated learning) cho bài toán \ac{ids} nhằm nâng cao khả năng thích ứng của mô hình và bảo vệ quyền riêng tư dữ liệu. Các phương pháp này cho phép phân tán quá trình huấn luyện trên nhiều nút hoặc thiết bị \ac{iot} khác nhau. Tuy nhiên, dữ liệu vẫn chủ yếu được xử lý theo từng đợt, thiếu cơ chế xử lý luồng tập trung và chưa khai thác đầy đủ khả năng mở rộng cũng như song song hóa của các nền tảng Big Data \cite{baidar2025,williams2025,albanbay2025,nguyen2019diot,boukela2021}.

Từ tổng quan các công trình nghiên cứu nêu trên, có thể nhận thấy rằng phần lớn các hệ thống \ac{ids} hiện tại vẫn hoạt động trong môi trường ngoại tuyến, chưa hỗ trợ xử lý dữ liệu liên tục theo thời gian thực và chưa tận dụng được khả năng mở rộng của các nền tảng xử lý dữ liệu lớn. Do đó, luận văn này hướng đến việc tích hợp các mô hình học máy dựa trên nền tảng Apache Spark nhằm xây dựng một hệ thống phát hiện xâm nhập có khả năng xử lý dữ liệu lớn, hỗ trợ thời gian thực và áp dụng hiệu quả trong môi trường mạng \ac{iot} hiện đại.

\section{Mục tiêu, đối tượng và phạm vi nghiên cứu}
\subsection{Mục tiêu nghiên cứu}
\label{sec:MucTieu}
Mục tiêu của đề tài “PHÁT HIỆN XÂM NHẬP (\ac{ids}) DỰA TRÊN APACHE SPARK CHO BẢO MẬT MẠNG IOT” là xây dựng một hệ thống phát hiện xâm nhập có khả năng phân tích và phát hiện các hành vi bất thường hoặc tấn công mạng trong môi trường \ac{iot} bằng cách kết hợp các mô hình học máy với nền tảng xử lý dữ liệu lớn Apache Spark nhằm đạt được:
\begin{itemize}
  \item Phát hiện hiệu quả các dạng tấn công mạng phổ biến (như DoS, probing, brute-force) trên các tập dữ liệu hiện đại như CIC-\ac{ids}2017.
  \item Thiết kế và triển khai pipeline xử lý dữ liệu phân tán và song song sử dụng Apache Spark, bao gồm các bước: tiền xử lý dữ liệu, trích xuất đặc trưng, huấn luyện mô hình và phát hiện thời gian thực.
  \item Tối ưu hóa hiệu suất hệ thống, đảm bảo khả năng mở rộng, độ trễ thấp và tính linh hoạt khi triển khai trên các hệ thống lớn hoặc mạng \ac{iot} nhiều thiết bị.
  \item So sánh và đánh giá các mô hình học máy phổ biến như Decision Tree, Support Vector Machine, Logistic Regression, \ac{nb}, Random Forest, \ac{gbt}, \ac{xgb} và các thuật toán Ensemble learning trong môi trường xử lý phân tán, từ đó chọn ra mô hình phù hợp nhất.
  \item Áp dụng quy trình đánh giá nghiêm ngặt, loại trừ rò rỉ dữ liệu (loại đặc trưng rò rỉ nhãn, khử trùng lặp ở mức đặc trưng, kiểm định ý nghĩa thống kê) nhằm đảm bảo kết quả phản ánh đúng năng lực mô hình.
\item Đề xuất giải pháp thực tiễn giúp phát hiện tấn công nhanh chóng và chính xác trong các hệ thống mạng hiện đại có yêu cầu xử lý dữ liệu lớn và liên tục.
\end{itemize}

\subsection{Đối tượng nghiên cứu}
\label{sec:DoiTuong}

Trong đề tài này, chúng tôi tập trung nghiên cứu và triển khai phương pháp phát hiện bất thường trong lĩnh vực bảo mật mạng \ac{iot} bằng cách tận dụng sức mạnh của nền tảng xử lý dữ liệu lớn Apache Spark. Với sự phát triển nhanh chóng của các thiết bị \ac{iot}, vấn đề đảm bảo an toàn và bảo mật cho hệ thống ngày càng trở nên cấp thiết, đặc biệt là trong việc phát hiện các hành vi bất thường hoặc tấn công mạng tiềm ẩn.

Apache Spark, với khả năng xử lý dữ liệu phân tán và tốc độ tính toán cao, cho phép xử lý khối lượng lớn dữ liệu sinh ra từ các thiết bị \ac{iot} trong thời gian thực. Trong khuôn khổ nghiên cứu này, chúng tôi sử dụng các tính năng của Spark như xử lý dòng dữ liệu (streaming), học máy phân tán (MLlib) và phân tích dữ liệu quy mô lớn để xây dựng mô hình phát hiện bất thường hiệu quả. Qua đó, đề tài không chỉ đề xuất một giải pháp kỹ thuật có tính ứng dụng cao mà còn góp phần nâng cao hiệu quả bảo mật cho các hệ thống \ac{iot} trong thực tiễn.

\subsection{Phạm vi nghiên cứu}
\label{sec:PhamVi}
Đề tài tập trung vào việc nghiên cứu và áp dụng các kỹ thuật phân tích dữ liệu lớn nhằm phát hiện bất thường trong hệ thống mạng \ac{iot} với mục đích là khai thác những điểm mạnh của nền tảng Apache Spark để xử lý và phân tích dữ liệu hiệu quả. Cụ thể, phạm vi nghiên cứu bao gồm:
\begin{itemize}
  \item \textbf{Dữ liệu đầu vào}: Dữ liệu được sử dụng là tập dữ liệu công khai CIC-\ac{ids}2017.
  \item \textbf{Công nghệ sử dụng}: Sử dụng Apache Spark làm nền tảng xử lý dữ liệu lớn, đặc biệt tập trung vào các thành phần như Spark Core, Spark SQL, Spark Streaming và Spark MLlib để hỗ trợ tiền xử lý dữ liệu, huấn luyện mô hình học máy và phát hiện bất thường theo thời gian thực.
  \item \textbf{Phương pháp phát hiện bất thường}: Áp dụng các thuật toán học máy được triển khai và tối ưu hóa trong môi trường Spark để đảm bảo khả năng mở rộng và hiệu suất cao.
  \item \textbf{Môi trường triển khai}: Thử nghiệm hệ thống trên một môi trường giả lập mạng \ac{iot} và bộ dữ liệu mô phỏng công khai (CIC-\ac{ids}2017) nhằm đánh giá khả năng phát hiện bất thường.
\end{itemize}

\section{Phương pháp nghiên cứu}
\label{sec:PhuongPhap}

\begin{itemize}
  \item \textbf{Thu thập, phân tích, xử lý dữ liệu}: Sử dụng tập dữ liệu CIC-\ac{ids}2017, từ đó tiến hành phân tích, xử lý và trực quan hoá dữ liệu.
  \item \textbf{Nghiên cứu về phương pháp tối ưu hoá siêu tham số, các phương pháp giảm chiều dữ liệu}: Nghiên cứu các phương pháp giảm chiều dữ liệu như \ac{pca}, \ac{shap} và Random Forest Importance nhằm giảm nhiễu và tăng độ chính xác, đồng thời tối ưu hóa siêu tham số thông qua Grid Search và Cross-Validation giúp nâng cao hiệu suất mô hình. Các bước này được triển khai trên nền tảng Apache Spark nhằm đảm bảo khả năng xử lý dữ liệu lớn hiệu quả và mở rộng.
  \item \textbf{Nghiên cứu và đề xuất mô hình học máy, học sâu, học máy tổng hợp phát hiện bất thường sử dụng Spark Machine Learning}: Sử dụng nhiều thuật toán khác nhau như Decision Tree, Support Vector Machine, Logistic Regression, \ac{nb}, Random Forest, \ac{gbt}, \ac{xgb} và các thuật toán Ensemble learning để xây dựng mô hình nhằm so sánh, đánh giá mô hình có hiệu suất tốt nhất.
  \item \textbf{Thực nghiệm và đánh giá hiệu năng mô hình trên cụm Jetson Orin Nano Super}: Xây dựng một hệ thống phát hiện xâm nhập mạng \ac{ids} \emph{phân tán} triển khai trên cụm hai thiết bị biên Jetson Orin Nano Super (cổng phát hiện bất thường trên một nút, bộ phân loại PySpark trên nút còn lại) nhằm đánh giá khả năng hoạt động trong môi trường tài nguyên hạn chế. Hệ thống được thiết kế theo kiến trúc xử lý dữ liệu thời gian thực, tận dụng Apache Kafka cho cơ chế truyền tải thông điệp, Spark Streaming cho xử lý luồng dữ liệu, PostgreSQL để quản lý dữ liệu, InfluxDB để lưu trữ chỉ số hiệu năng theo thời gian, Grafana để trực quan hóa và Mailtrap phục vụ kiểm thử cơ chế cảnh báo qua email.
\end{itemize}

\section{Các đóng góp chính của đề tài}
\label{sec:DongGop}
\begin{itemize}
  \item \textbf{Đề xuất giải pháp ứng dụng Apache Spark trong phát hiện bất thường cho hệ thống mạng \ac{iot}}: Đề tài đã xây dựng một mô hình phát hiện bất thường dựa trên nền tảng xử lý dữ liệu lớn Apache Spark, tận dụng khả năng xử lý phân tán và thời gian thực của Spark để đáp ứng yêu cầu về hiệu suất và tính mở rộng trong môi trường \ac{iot} có khối lượng dữ liệu lớn và biến đổi liên tục.
  \item \textbf{Xây dựng quy trình xử lý và phát hiện bất thường toàn diện trên dữ liệu \ac{iot}}: Từ bước tiền xử lý dữ liệu, trích xuất đặc trưng, huấn luyện mô hình đến đánh giá kết quả, đề tài đã thiết kế và hiện thực một quy trình phát hiện bất thường hoàn chỉnh có thể áp dụng cho các hệ thống thực tế.
  \item \textbf{ Ứng dụng và đánh giá các thuật toán học máy trong môi trường phân tán}: Đề tài triển khai và so sánh hiệu quả của một số thuật toán học máy và học sâu trên nền tảng Apache Spark MLlib, qua đó làm rõ ưu điểm và hạn chế của từng mô hình trong việc xử lý dữ liệu \ac{iot}.
  \item \textbf{Đánh giá hiệu năng hệ thống trong môi trường dữ liệu lớn}: Không chỉ đánh giá độ chính xác của mô hình, đề tài còn phân tích hiệu năng xử lý (thời gian thực thi, khả năng mở rộng) của hệ thống trên các tập dữ liệu lớn nhằm chứng minh tính thực tiễn của giải pháp trong môi trường triển khai thực tế.
  \item \textbf{Quy trình đánh giá loại trừ rò rỉ dữ liệu và có kiểm định thống kê}: Đề tài nhận diện và loại bỏ các nguồn rò rỉ dữ liệu thường gặp trên CIC-\ac{ids}2017 (đặc trưng \texttt{destination\_port} rò rỉ nhãn, trùng lặp luồng ở mức đặc trưng), đồng thời áp dụng kiểm định ý nghĩa thống kê, giúp kết quả phản ánh đúng năng lực mô hình thay vì các con số bị thổi phồng.
\end{itemize}

\section{Công bố khoa học}
Một phần kết quả của luận văn được phát triển thành hai công trình khoa học (đồng tác giả Nguyễn Vũ Thái và Bùi Văn Dũng), hiện đang được hoàn thiện để gửi đăng:
\begin{enumerate}
  \item Nguyễn Vũ Thái, Bùi Văn Dũng, ``So sánh giảm chiều và lựa chọn đặc trưng cho phát hiện xâm nhập mạng trên Apache Spark ML với quy trình đánh giá loại trừ rò rỉ dữ liệu,'' Hội nghị Nghiên cứu cơ bản và ứng dụng Công nghệ thông tin (FAIR), 2026.
  \item Nguyen Vu Thai, Bui Van Dung, ``A Distributed Real-Time Intrusion Detection System on a Jetson Orin Nano Super Edge Cluster using Kafka and PySpark,'' Symposium on Information and Communication Technology (SOICT), 2026.
\end{enumerate}

\section{Cấu trúc của luận văn}
\label{sec:CauTruc}
Khóa luận với đề tài ``PHÁT HIỆN XÂM NHẬP (\ac{ids}) DỰA TRÊN APACHE
SPARK CHO BẢO MẬT MẠNG IOT'' được trình bày bao gồm 5 chương. Nội dung tóm tắt từng chương được trình bày như sau:
\begin{itemize}
    \item \textbf{Chương 1: Mở đầu}: Trình bày tổng quan về bối cảnh và lý do chọn đề tài, mục tiêu nghiên cứu, phạm vi nghiên cứu, câu hỏi nghiên cứu và phương pháp tiếp cận. Ngoài ra, chương này cũng nêu rõ ý nghĩa thực tiễn và khoa học của đề tài, cũng như bố cục tổng thể của luận văn.
    \item \textbf{Chương 2: Cơ sở lý thuyết}: Trình bày các khái niệm cơ bản và nền tảng lý thuyết phục vụ cho đề tài, bao gồm tổng quan về mạng \ac{iot}, các lỗ hổng và mối đe dọa bảo mật, phương pháp phát hiện xâm nhập (\ac{ids}), và kiến thức liên quan đến Apache Spark. Ngoài ra, chương này cũng phân tích các thuật toán học máy thường được sử dụng để phát hiện bất thường.
    \item \textbf{Chương 3: Phương pháp thực hiện}: Trình bày chi tiết các bước xây dựng hệ thống phát hiện xâm nhập, bao gồm thu thập và tiền xử lý dữ liệu, thiết kế kiến trúc hệ thống sử dụng Apache Spark, lựa chọn và triển khai các thuật toán học máy, cùng quy trình huấn luyện và đánh giá mô hình.
    \item \textbf{Chương 4: Thực nghiệm, đánh giá và thảo luận}: Trình bày kết quả thực nghiệm trên bộ dữ liệu thực tế hoặc mô phỏng, đánh giá hiệu suất của các mô hình đã triển khai theo các tiêu chí như độ chính xác, độ nhạy, F1-score, kích thước của mô hình và thời gian xử lý. Đồng thời, chương này cũng phân tích, so sánh và thảo luận về các kết quả đạt được nhằm làm rõ hiệu quả của phương pháp đề xuất.
    \item \textbf{Chương 5: Kết luận và hướng phát triển}: Tổng kết những kết quả chính mà đề tài đạt được, nêu bật những đóng góp chính, đồng thời chỉ ra những hạn chế còn tồn tại và đề xuất các hướng phát triển trong tương lai, như việc tích hợp các mô hình học sâu hoặc triển khai hệ thống trên nền tảng đám mây để nâng cao tính khả dụng và khả năng mở rộng.
\end{itemize}